\documentclass[11pt]{article}
\usepackage[margin=1in]{geometry}
\usepackage{amsmath,amssymb,amsthm,booktabs,xcolor}
\newtheorem{proposition}{Proposition}
\newtheorem{remark}{Remark}
\newcommand{\red}[1]{\textcolor{red}{#1}}
\newcommand{\code}[1]{\texttt{#1}}

\title{Space--Time--Behavior Flow-Through Tensors:\\
A Certified Multi-Axis Operator Framework Connecting\\
Logit SUE, Fluid Queues, and Layered Network Equilibria\\
{\large (STB-FTT v0.6 wrap-up draft --- 2026-07-20)}}
\author{Xuesong Zhou \and \red{[co-authors TBD]}}
\date{Draft v0.6 --- every number reproduced from the certified test
suite (49 tests) and run scripts of \code{stb\_ftt\_unified\_v0.6}}

\begin{document}
\maketitle

\begin{abstract}
The Flow-Through Tensor (FTT) architecture connects origin--destination
flows, path probabilities, and link travel times as chained tensors on
one computational graph. This paper wraps that foundation into a
\emph{certified multi-axis operator framework}: every axis of the
space--time--behavior tensor declares exactly a \emph{status}
(spectator, contracted, or synchronized), a \emph{semiring}
($(+,\times)$ flow, $(\min,+)$ pricing, $(\max,\times)$ state decode),
and a methodological \emph{anchor}; a model is a chain of promotions of
axes from spectator to synchronized; and every synchronized axis owes
one well-posedness argument and one certificate class. Three
foundational special cases close the framework from below --- logit
stochastic user equilibrium and its deterministic limit (behavior),
Newell fluid point queues (time), and the flow-through tensor chain
itself (space) --- each verified against closed forms inside one code
base. Above them, the same operator calculus carries multi-layer
passenger--freight price consensus, the passenger--vehicle service
coupling $Rx \le Sy$, and a mixed-autonomy ride-hailing profile with
competing companies, AV/HV fleets, and patience-capped matching queues.
A solver-engine escalation ladder (damped fixed point $\to$ averaging
$\to$ Anderson $\to$ stiff-block Newton $\to$ semismooth NCP/VI)
computes the fixed points that pin synchronized axes, with cycle
detection and engine bookkeeping as part of the certificate. All
claims are bounded by an explicit validity domain: learned or
behavioral terms are flow-independent along synchronized axes,
compression is admissible only on column-rich spectator axes, and
nonsmooth dynamic loading is a stated physics boundary, not a hidden
smoothing.
\end{abstract}

\section{Introduction}

Three research lines meet in this paper. First, the FTT computational
graph, which chains $\mathbf f_{OD} \to \mathbf f_P \to \mathbf f_L \to
\mathbf t_L \to \mathbf t_{OD}$ through typed incidence operators and
exploits the transpose symmetry of contraction for gradients. Second,
the classical equilibrium canon --- logit SUE, deterministic UE, and
fluid queue loading --- whose recovery as special cases is the ground
truth any unifying framework must pass. Third, layered service systems
--- multi-commodity freight, passenger--vehicle service coupling, and
mixed-autonomy ride-hailing --- whose defining feature is that layers
interact \emph{only} through shared resources synchronized by prices.

The claim of this paper is deliberately narrow, in the sense fixed by
our earlier position notes: we do not claim a new three-way unification
of neural, optimization, and transportation architectures. We claim a
\emph{certified operator framework}: (i) an axis calculus in which the
classical models are exact reductions, verified numerically to
closed-form accuracy; (ii) a scalable column-generation layer whose
optimality gaps are always priced over the full network, never over the
enumerated pool; (iii) a layered-equilibrium layer whose prices are
certified against exact duals where the coupling is convex, and whose
complementarity conditions are audited where it is not; and (iv) an
explicit engine ladder for the fixed points, with failure modes
diagnosed rather than tuned away.

\section{The axis calculus}

Let the demand tensor carry indices
$(o,d,g,m,\tau,p,a,t)$ --- origin, destination, group, mode, departure
period, path/column, link, entry time --- extended by profile-specific
axes (company, vehicle type, operational stage, matching queue). Every
axis declares exactly three properties:

\begin{center}
\begin{tabular}{@{}lll@{}}
\toprule
declaration & values & meaning \\
\midrule
status & spectator / contracted / synchronized &
rides along / summed out / pinned by a fixed point \\
semiring & $(+,\times)$, $(\min,+)$, $(\max,\times)$ &
flow aggregation / pricing / state decode \\
anchor & $S_{\mathrm{repr}}, S_{\mathrm{stat}}, S_{\mathrm{mech}},
S_{\mathrm{ctrl}}$ & which methodological layer owns it \\
\bottomrule
\end{tabular}
\end{center}

\noindent Four consequences follow mechanically.

\textbf{Models are promotion chains.} Static UE promotes only the link
axis (congestion fixed point). Dynamic assignment additionally promotes
departure/entry time (queue conservation recursion). Multi-layer
systems promote the layer axis (price consensus). State estimation
promotes the regime axis under $(\max,\times)$. The mixed-autonomy
profile promotes the matching-queue axis (patience-capped waiting).

\textbf{Compression lives on spectator axes only, and only where
columns are rich.} The Chicago Sketch experiment below shows the
boundary empirically: at $\bar K \approx 2.2$ active columns per OD,
route-family atoms cannot pay.

\textbf{Learning is admissible only if flow-independent along every
synchronized axis} (the clean-VI boundary): priors, values, and utility
offsets may be learned; anything entering a flow-dependent cost voids
the certificates unless monotone by construction.

\textbf{Each synchronized axis owes one certificate class:}
monotonicity $\Rightarrow$ FW/KKT gap; causality $\Rightarrow$
conservation and closed-form pulses; convex duality $\Rightarrow$
consensus gap against exact duals; integrality or nonmonotone coupling
$\Rightarrow$ report the residual gap, never assume it away.

\section{Special-case closure}

All three foundational cases are certified inside one package
(\code{tests/test\_v06\_closure.py} and the Sioux Falls special-case
suite; data: TCGlite GMNS Sioux Falls, 24 nodes / 76 links / 33 ODs /
126 behavioral paths).

\subsection{Behavior: logit SUE and the deterministic limit}

The entropy-regularized master over the fixed path pool has, at frozen
costs, the multinomial logit as its \emph{exact} KKT solution (machine
precision, $0.0$ error). With congestion, Evans-type partial
linearization with exact line search yields self-consistency residuals
$\le 2\times10^{-5}$ across $\theta \in [0.1, 50]$, and the Wardrop gap
falls monotonically from $6.9\times10^{-2}$ at $\tau{=}10$ to
$8.3\times10^{-6}$ at $\tau{=}0.02$: the $\tau \to 0$ ladder reaches
deterministic UE inside one solver family. Two solver lessons are part
of the record: the plain damped fixed point \emph{oscillates} at high
congestion (residual $0.99$), and demand scaling matters ($\times 25$
gives $v/c = 4.7$ and an ill-conditioned master; $\times 8$, $v/c
\approx 2$, is the certified stress level).

A full-network Dijkstra audit prices the fixed 126-path pool at every
reported solution: the pool is \emph{exact} at free flow (relative gap
$1.9\times10^{-16}$) and \emph{35.7\% inadequate} at congested UE
times. The static-pool failure mode is therefore measured, not
assumed, and adaptive column generation is the certified remedy.

\subsection{Time: Newell fluid point queues}

The point-queue recursion reproduces the closed-form constant-overload
pulse ($\lambda{=}30$, $\mu{=}20$, $\tau_0{=}10$ min) to relative
errors $\le 7\times10^{-5}$ in peak queue, clearing time, and total
delay; network flow-through of the certified UE column flows conserves
mass to $9\times10^{-12}$ vehicles over 71 loaded links, with causality
$D(t)\le A(t)$ everywhere. The boundary is stated with the operator:
sending/receiving $\min(\cdot)$ is nonsmooth at binding capacity;
smoothing restores gradients but changes the physics; no
differentiability through the queue is claimed.

\subsection{Space: the flow-through tensor chain}

The TCG chain $\alpha \to \Gamma \to P \to \Delta$ (zone generation,
OD split, path proportion, link incidence) is exactly the conserved
decoder composition $x = D(\Pi)X$: row-stochastic masked layers, mass
conservation $2\times10^{-13}$, link reconstruction residual
$<10^{-9}$, and the forward chain at target observations reproduces the
sensor counts to the data's own rounding (max relative error
$1.0\times10^{-2}$). The calibration direction carries an analytic
gradient through the masked-softmax chain, verified by finite
differences to $2\times10^{-7}$ --- with the methodological caveat that
the first FD check was vacuous (masked coordinates have identically
zero gradient and difference) until sampling was restricted to
meaningful coordinates.

\section{Scale and the compression boundary}

The same solvers run unchanged on a deterministic grid generator, Sioux
Falls, and Chicago Sketch (933 nodes, 2{,}950 links, 386 zones,
93{,}135 ODs, 1.14M trips) through one canonical GMNS contract.
Full-network certified results (single process, single thread; all
gaps priced by all-origin Dijkstra over the real network):

\begin{center}
\begin{tabular}{@{}lrrrr@{}}
\toprule
algorithm & full-space rel.\ gap & wall (s) & columns & atoms \\
\midrule
CG Frank--Wolfe & $8.3\times10^{-5}$ & 7.4 & 198{,}871 & --- \\
CG FW + grouped-simplex GP & $9.4\times10^{-5}$ & 39.3 & 207{,}248 & --- \\
route-family atoms (static) & $6.1\times10^{-3}$ & 141.9 & 200{,}712 & 157{,}373 \\
route-family atoms (adaptive) & $4.6\times10^{-3}$ & 178.4 & 215{,}237 & 220{,}152 \\
\bottomrule
\end{tabular}
\end{center}

\noindent The certified link flows correlate $1.0000$ with the
network's reference volumes (RMSE 23.8 veh). The atom rows are the
honest boundary datum: at $\bar K \approx 2.2$ columns per OD, the
compression ratio is only $1.28\times$, adaptive promotion generates
62{,}761 singletons until atoms \emph{outnumber} columns, and the gap
floor sits two orders above tolerance --- reproducing, at scale, the
crossover law that compression pays only in column-rich regimes
(activity hierarchies, transit itineraries), never in the coupling
itself. Wall-clock numbers support algorithm ordering only; engine-level
speed claims are out of scope.

\section{Layers and prices}

\subsection{Convex two-layer benchmark with exact dual certificates}

Passenger columns (mandatory, PCE 1) and elective e-commerce freight
requests (PCE 2) share link capacity on Sioux Falls. With linear costs
the joint problem is an LP: 11 binding links, max price 40.8
min/veh-equiv. Iterative block coordination --- parallel min-plus best
responses per layer, projected subgradient on the shared price, primal
averaging --- certifies against the exact optimum: dual-bound gap
$3.3\times10^{-4}$, price agreement within 1.6\% on every binding
link, served freight $38.1\%$ falling to $32.4\%$ as slack tightens.
Replacing hard capacity by BPR congestion keeps joint convexity (one
Beckmann potential in total veh-equivalent load): joint FW gap
$1.8\times10^{-6}$; Gauss--Seidel block coordination agrees to
$5.3\times10^{-7}$ with link loads within 0.12 veh-equiv; freight
threshold complementarity is exact; and congestion \emph{alone} prices
freight from 100\% served ($\times4$ passenger demand) to 86.3\%
($\times16$) --- the congested travel time is itself the synchronizing
price. The step-size lesson is recorded: the subgradient step must be
scaled to the load/price ratio ($\rho_0{=}0.5$ diverges; $0.01$
converges; report the best dual bound).

\subsection{The passenger--vehicle coupling and the integer corner}

The service coupling $Rx \le Sy$ with a scarcity price on service
capacity is carried in the package as its own layer, with passenger
mass, service capacity, vehicle flow, and queue mass audited
separately. The activity--travel vehicle-scheduling line (the
``how many trip requests'' formulation) is precisely the \emph{integer
corner} of this master: passenger chains as columns, one-chain
constraints as simplexes, on-board arc coupling as incidence, and the
Lagrangian trip price as the synchronizing dual. On the integer side
the duality gap need not vanish and must be reported; our convex
benchmarks are the certified relaxation of that corner.

\subsection{The mixed-autonomy ride-hailing profile}

The MAGE class (competing TNCs; AV/HV fleets under an AV-share cap;
pickup and service stages; patience-capped customer matching;
congestion-coupled routing) enters as an \emph{operator profile}:
extension axes (company, vehicle type, stage, matching queue) are
registered without touching the canonical registry; the matching queue
is kept structurally distinct from the road queue; and the road layer
reuses the certified full-space assignment each outer iteration.
On an $8\times8$ grid with two asymmetric companies the equilibrium
closes in 4 outer iterations (inner residual $4\times10^{-14}$;
conservation $10^{-13}$; patience, rationing-complementarity, and
fleet-cap violations $0$; multi-start share spread $4\times10^{-6}$).
Comparative statics behave economically: raising one company's AV cap
from $0$ to $1$ moves its AV share from 6.7\% to 29.8\% and cuts its
matching wait from 12.0 to 1.8 minutes; fleet scale $\times0.3 \to
\times1.5$ moves served fraction from 46\% to 100\%; and the patience
sweep exposes the defining subtlety --- \emph{lower} patience (4 min)
\emph{raises} ride-hailing share to 64.8\% while service is rationed to
77.9\%, because patience caps the disutility, not the supply, with the
complementarity (unserved ${>}\,0 \Rightarrow$ wait at patience)
holding exactly. The full NCP/VI equilibrium (PATH class) is the
declared next solver rung; the profile's existence claims are
restricted to the steady-state operators, exactly as in the source
formulation.

\section{Equilibrium engines}

Fixed points that pin synchronized axes are computed by an explicit
escalation ladder --- damped Picard (E0), MSA averaging (E1),
safeguarded Anderson acceleration (E2), stiff-block Newton (E3), and
semismooth NCP/VI (E4, declared) --- with a period-$p$ cycle detector
as the promotion trigger and the engine choice recorded in the result.
The motivating instance is part of the test suite as an honest
negative: the ride-hailing choice$\leftrightarrow$wait subsystem has
$|\partial w / \partial s| \sim 10^4$ near saturation, damped iteration
limit-cycles at residual $0.12$ (detected, period 4), smoothing and
two-timescale damping do not rescue it, and the exact Newton solve of
the \emph{smallest stiffest block in its low-dimensional dual-like
variable} (the wait vector) closes it. That design rule --- solve the
stiff block exactly in waits/prices/multipliers, never in
high-dimensional primal shares; keep slow blocks in cheap fixed-point
form --- is the solver-side counterpart of the
coordination-vs-synchronization distinction.

\section{Roadmap: the phase--time axis (v0.7)}

The deferred state/regime axis has a concrete target: the
phase-augmented DTA formulation with X (queue dynamics), Y (generalized
phase selection), and Z (vehicle trajectory) blocks coupled by
Lagrangian relaxation. In the axis calculus this is one further
promotion: the regime axis becomes synchronized, contracted under
$(\max,\times)$ for decoding and $(\min,+)$ for pricing on the
phase-expanded network, with within-phase closed-form queue profiles
replacing the binding $\min(\cdot)$ --- moving the nonsmoothness to the
phase boundary where a discrete engine (E4/branch-and-price) owns it.
The engine ladder and the consensus machinery above are exactly the
solver substrate this promotion requires.

\section{Validity domain and limitations}

(i) All learned/behavioral terms in every certified experiment are
flow-independent along synchronized axes; the monotonicity boundary is
stated wherever a flow-dependent extension would cross it.
(ii) The dynamic loading operator is forward-only; no gradients are
claimed through binding capacities.
(iii) Integer couplings (vehicle tours, phase selection) carry duality
gaps that are reported, not assumed zero.
(iv) The well-posedness map currently spans behavioral-feedback and
queue-feedback strength; the coupling-integrality axis is documented
but not yet swept.
(v) Wall-clock comparisons support algorithm ordering on one machine,
single-threaded; engine-level speed claims require the separate
compiled-solver line.

\section{Conclusion}

The FTT chain supplies the spine; the axis calculus supplies the
grammar; the special cases nail the framework to closed forms; the
layered benchmarks show that one price mechanism --- explicit duals
under hard coupling, congested times under soft coupling, patience-capped
waits under service coupling --- synchronizes every layer pair; and the
engine ladder makes the computation of those fixed points a diagnosed,
certified choice rather than a tuning exercise. Everything above is
reproducible from one package with one test suite (49 tests) and the
run scripts cited throughout.

\paragraph{Reproducibility.} \code{python -m pytest} (49 tests);
\code{run\_mage\_grid.py}, \code{run\_unified\_grid\_harness.py},
\code{run\_passenger\_vehicle\_harness.py} in
\code{stb\_ftt\_unified\_v0.6}; multi-layer benchmarks in
\code{STB\_multilayer\_benchmark} (7 tests); special-case artifacts in
\code{STB\_special\_cases\_SF} (7 tests). \red{[Bibliography, author
list, and venue formatting pending author decisions; FTT manuscript and
MAGE submission to be cited as companion/related work.]}

\end{document}
